\begin{claim}{E2.3}{Open}{The remaining comb splitting and deformation calculations}
The integer $m$ and the stated limiting data do not determine in
this paper the splitting of the rank-three bundle $K_m$ in (E.6),
its exact number of sections, or the obstruction map of the comb
member. These parts of E.2 are unfinished. The four tangent
cohomology dimensions themselves are fully determined by (E.4).
\end{claim}
\begin{proof}[The remaining finite data and the computations made]
If $e_m$ is the length of $\operatorname{coker}\kappa_1$ in (E.6),
then its source has rank four and degree $-2$, its target is
$\OO_B(2m)$, and
\[
 \deg K_m=-2m-2+e_m.                                  \tag{E.7}
\]
The degree $-2$ follows from the two pairs of opposite stabilizer
characters of $E_m^*\otimes G_m$; its root determinant is trivial,
and each finite root point subtracts one from the coarse determinant.
Although (E.4a) determines the tangent dimensions independently of
the individual zeros of the Kodaira--Spencer determinant, the
support and elementary divisors of $Q_m$ still depend on that
period-section map. The comb Kuranishi problem further requires
its actual bracket and higher obstruction maps; the dimensions in
(E.4) do not determine those maps.

We evaluated the universal exact degree and dimension expressions
for $m=12,13,20,100$ in \texttt{d\_exact.py} and
\texttt{d\_comb.py}; the first case already
gives $h^{1,1}\ge49$, $h^{1,2}\ge48$, $h^0(T)=25$,
$h^1(T)=97$ and $h^2(T)=73$. The smallest undecided instance
is the splitting of (E.6) for one specified smoothed 12-tooth period
section, together with that member's obstruction map. The
attached construction supplies an existence argument and an open
set of smoothing choices, not the transition matrices of a uniquely
specified $X_{12}$. We have not chosen and solved its nonlinear
period-section smoothing equations, and have not run two numerical
methods on those missing matrices. Formulas (E.5)--(E.7) and the
growth bounds are proved, but they are not replacements for that
remaining calculation or for its Kuranishi bracket.
\end{proof}

